\noindent\textbf{Elección del procedimiento.} Con gradiente $-0.0065$ K/m, $T=T_0-Gz$. Integrar hidrostática con gas ideal produce $p=p_0(T/T_0)^{g_0/(RG)}$. Después usar $\rho=p/(RT)$ conservando la presión sin redondear.

\begin{enumerate}

\item La temperatura disminuye linealmente en esta capa.

\[\fitmath{T=T_{0}-G\cdot z}\]

\[\fitmath{T=(288.15)-(0.0065)\cdot (5000)\approx 255.65\,\mathrm{K}}\]

\item Introducir el cociente de temperaturas en la ley barométrica.

\[\fitmath{p=p_{0}\cdot {\left(\frac{T}{T_{0}}\right)}^{\frac{g_{0}}{\mathrm{Ra}\cdot G}}}\]

\[\fitmath{p=(101325)\cdot {\left(\frac{(255.65)}{(288.15)}\right)}^{\frac{(9.80665)}{(287.05287)\cdot (0.0065)}}\approx 54019.89\,\mathrm{Pa}}\]

\item Cerrar el estado con la temperatura local.

\[\fitmath{\rho=\frac{p}{\mathrm{Ra}\cdot T}}\]

\[\fitmath{\rho=\frac{(54019.88819)}{(287.05287)\cdot (255.65)}\approx 0.7361155\,\mathrm{kg/m^{3}}}\]

\end{enumerate}
